% -------------------------------
% Plantilla del TFG en la ESIIAB-UCLM. Versión 0.0 (Preliminar)
%
% Luis de la Ossa
% 
% Compilar con XeLaTeX 
% -------------------------------
\documentclass[12pt, a4paper]{book}

% -------------------------------
% Información y opciones del documento. Debe editarse.
% -------------------------------
\input{include/opciones}

% -------------------------------
% Configuración del documento, paquetes, etc. 
% -------------------------------
\input{include/configuracion}

% Documento
\begin{document}

% -------------------------------
% Portada y título
% -------------------------------
% Fuente  de la portada
% \setmainfont[Ligatures={NoRequired,NoCommon,NoContextual}]{Calibri}
\input{elements/portada}

% -------------------------------
% Fuente del documento.
% -------------------------------
% La fuente recomendada es Calibri, pero no está instalada en overleaf:
%\setmainfont[Ligatures={NoRequired,NoCommon,NoContextual}]{Calibri}

% Para utilizarla aquí, son necesarios los archivos .ttf correspondientes:
%\setmainfont[
%			BoldFont=Calibrib.ttf,
%			ItalicFont=Calibrii.ttf,
%			BoldItalicFont=Calibriz.ttf
%		]{Calibri.ttf}

% Por defecto, utiliza una fuente instalada en overleaf. 
% Éstas son algunas que van bien con la plantilla:
%\setmainfont{Alegreya Sans}
\setmainfont{Carlito}
%\setmainfont{PT Sans}
%\setmainfont{Source Sans Pro}

% -------------------------------
% Preámbulo
% -------------------------------
\frontmatter 
\input{elements/preambulo}

% -------------------------------
% Índices y tablas
% -------------------------------
\tableofcontents	% Índice
\renewcommand{\listfigurename}{List of Figures} % Índice de figuras
\listoffigures		
\renewcommand{\listtablename}{List of Tables} % Índice de tablas
\listoftables		
\listofalgorithms   % Índice de algoritmos (comentar si no los hay)
\listofcodes	    % Índice de códigos (comentar si no los hay)
\cleardoublepage

% -------------------------------
% Documento
% -------------------------------
\mainmatter
% Estilo de página
\pagestyle{fancy}

% Introducción
\input{chapters/ch1} 

% Algunos elementos
\input{chapters/ch2} 

% Utilidades para anotaciones
\input{chapters/ch3} 

% -------------------------------
% Apéndices
% -------------------------------
\appendix
\chapter{Annex 1}
\lipsum[1-2]

\cleardoublepage
\thispagestyle{empty}


% -------------------------------
% Bibliografía
% -------------------------------
\bibliographystyle{apalike}
\bibliography{bib/ref} 
\addcontentsline{toc}{chapter}{Bibliography} 

\end{document}